\documentclass[11pt]{article}
\newcommand{\LabNumber}{4.1}
\newcommand{\LabTitle}{Switches, LEDs, 7-Seg Mux}
\newcommand{\LabTopic}{Breadboard I/O shield for the Arduino Uno R3 (push buttons, LEDs, and a multiplexed 7-segment display), tested with AVR assembly}
\newcommand{\LabDuration}{3 hours (in-lab)}
% Shared preamble for Computer Programming lab sheets.
% Include with % Shared preamble for Computer Programming lab sheets.
% Include with % Shared preamble for Computer Programming lab sheets.
% Include with \input{common.tex} after \documentclass{article}.

\usepackage[a4paper,margin=2.2cm]{geometry}
\usepackage[T1]{fontenc}
\usepackage{lmodern}
\usepackage{microtype}
\usepackage{enumitem}
\usepackage{xcolor}
\usepackage{tcolorbox}
\tcbuselibrary{breakable,skins}
\usepackage{listings}
\usepackage{fancyhdr}
\usepackage{titlesec}
\usepackage{amsmath}
\usepackage{array}
\usepackage{booktabs}
\usepackage{hyperref}

\hypersetup{
  colorlinks=true,
  linkcolor=black,
  urlcolor=blue!60!black,
  pdfborder={0 0 0}
}

% ---------- Course identity (override per file if needed) ----------
\providecommand{\CourseName}{Microprocessors}
\providecommand{\CourseCode}{010153521}
\providecommand{\Semester}{Semester 1/2026}
\providecommand{\LabDuration}{3 hours (in-lab)}

% ---------- Colors ----------
\definecolor{accent}{HTML}{1F4E79}
\definecolor{accentlight}{HTML}{D9E2EC}
\definecolor{warn}{HTML}{B45309}
\definecolor{ok}{HTML}{166534}
\definecolor{codebg}{HTML}{F6F8FA}
\definecolor{codeframe}{HTML}{D0D7DE}
\definecolor{codekw}{HTML}{0550AE}
\definecolor{codestr}{HTML}{0A7B2C}
\definecolor{codecom}{HTML}{6E7781}

% ---------- Code / assembly listings ----------
\lstdefinestyle{c}{
  language=C,
  basicstyle=\ttfamily\small,
  keywordstyle=\color{codekw}\bfseries,
  stringstyle=\color{codestr},
  commentstyle=\color{codecom}\itshape,
  backgroundcolor=\color{codebg},
  frame=single,
  rulecolor=\color{codeframe},
  framerule=0.4pt,
  showstringspaces=false,
  breaklines=true,
  columns=fullflexible,
  keepspaces=true,
  tabsize=4,
  xleftmargin=8pt,
  xrightmargin=4pt,
  aboveskip=6pt,
  belowskip=6pt,
}
\lstdefinestyle{asm}{
  basicstyle=\ttfamily\small,
  commentstyle=\color{codecom}\itshape,
  backgroundcolor=\color{codebg},
  frame=single,
  rulecolor=\color{codeframe},
  framerule=0.4pt,
  showstringspaces=false,
  breaklines=true,
  columns=fullflexible,
  keepspaces=true,
  tabsize=4,
  xleftmargin=8pt,
  xrightmargin=4pt,
  aboveskip=6pt,
  belowskip=6pt,
}
\lstset{style=asm}

% ---------- Section styling ----------
\titleformat{\section}{\Large\bfseries\color{accent}}{\thesection.}{0.5em}{}
\titleformat{\subsection}{\large\bfseries\color{accent!85!black}}{\thesubsection}{0.5em}{}
\titlespacing*{\section}{0pt}{12pt}{6pt}
\titlespacing*{\subsection}{0pt}{8pt}{4pt}

% ---------- Page headers/footers ----------
\pagestyle{fancy}
\fancyhf{}
\fancyhead[L]{\small\CourseName\ifx\CourseCode\empty\else\ (\CourseCode)\fi}
\fancyhead[C]{\small\Semester}
\fancyhead[R]{\small\textbf{Lab \LabNumber:} \LabTitle}
\fancyfoot[C]{\small\thepage}
\renewcommand{\headrulewidth}{0.4pt}

% ---------- Title block ----------
\newcommand{\labheader}{%
  \begin{tcolorbox}[
    enhanced,
    colback=accentlight,
    colframe=accent,
    boxrule=0.5pt,
    arc=2pt,
    left=10pt, right=10pt, top=8pt, bottom=8pt
  ]
    {\Large\bfseries\color{accent} Lab \LabNumber: \LabTitle}\\[2pt]
    {\small \CourseName\ifx\CourseCode\empty\else\ (\CourseCode)\fi\quad\textbar\quad \Semester}\\[1pt]
    {\small \textbf{Duration:} \LabDuration\quad\textbar\quad \textbf{Topic:} \LabTopic}
  \end{tcolorbox}
  \vspace{2pt}
}

% ---------- Custom environments ----------
\newtcolorbox{summary}[1][Quick Reference]{
  enhanced,breakable,
  colback=accentlight!40,colframe=accent!70,
  title={\bfseries #1},
  fonttitle=\color{white},
  attach boxed title to top left={xshift=6pt,yshift=-3pt},
  boxed title style={colback=accent,boxrule=0pt,arc=2pt},
  arc=2pt,boxrule=0.4pt,
  top=12pt,
}

\newcounter{labex}[section]
\newtcolorbox{labexercise}[1]{
  enhanced,breakable,
  colback=white,colframe=accent,
  title={\bfseries In-Lab Exercise \refstepcounter{labex}\thelabex: #1},
  fonttitle=\color{white},
  attach boxed title to top left={xshift=6pt,yshift=-3pt},
  boxed title style={colback=accent,boxrule=0pt,arc=2pt},
  arc=2pt,boxrule=0.4pt,
  top=12pt,
}

\newcounter{traceex}
\newtcolorbox{trace}[1][]{
  enhanced,breakable,
  colback=white,colframe=warn,
  title={\bfseries Trace \& Predict \refstepcounter{traceex}\thetraceex\ifx\\#1\\\else: #1\fi},
  fonttitle=\color{white},
  attach boxed title to top left={xshift=6pt,yshift=-3pt},
  boxed title style={colback=warn,boxrule=0pt,arc=2pt},
  arc=2pt,boxrule=0.4pt,
  top=12pt,
}

\newcounter{bugex}
\newtcolorbox{bug}[1][]{
  enhanced,breakable,
  colback=white,colframe=warn,
  title={\bfseries Find the Bug \refstepcounter{bugex}\thebugex\ifx\\#1\\\else: #1\fi},
  fonttitle=\color{white},
  attach boxed title to top left={xshift=6pt,yshift=-3pt},
  boxed title style={colback=warn,boxrule=0pt,arc=2pt},
  arc=2pt,boxrule=0.4pt,
  top=12pt,
}

\newcounter{writeex}
\newtcolorbox{writecode}[1][]{
  enhanced,breakable,
  colback=white,colframe=warn,
  title={\bfseries Write Code on Paper \refstepcounter{writeex}\thewriteex\ifx\\#1\\\else: #1\fi},
  fonttitle=\color{white},
  attach boxed title to top left={xshift=6pt,yshift=-3pt},
  boxed title style={colback=warn,boxrule=0pt,arc=2pt},
  arc=2pt,boxrule=0.4pt,
  top=12pt,
}

% ---------- AI Collaboration box ----------
\definecolor{aipurple}{HTML}{4C1D95}
\definecolor{aipurplelight}{HTML}{EDE9FE}

\newtcolorbox{aicollab}[1][AI Collaboration]{
  enhanced,breakable,
  colback=aipurplelight,colframe=aipurple!70,
  title={\bfseries #1},
  fonttitle=\color{white},
  attach boxed title to top left={xshift=6pt,yshift=-3pt},
  boxed title style={colback=aipurple,boxrule=0pt,arc=2pt},
  arc=2pt,boxrule=0.4pt,
  top=12pt,
}

\newcounter{tryex}
\newtcolorbox{tryit}[1]{
  enhanced,breakable,
  colback=ok!6,colframe=ok!70!black,
  title={\bfseries Try It \refstepcounter{tryex}\thetryex: #1},
  fonttitle=\color{white},
  attach boxed title to top left={xshift=6pt,yshift=-3pt},
  boxed title style={colback=ok!70!black,boxrule=0pt,arc=2pt},
  arc=2pt,boxrule=0.4pt,
  top=12pt,
}

% ---------- Drawing exercise ----------
\definecolor{drawteal}{HTML}{0D7377}
\definecolor{drawteallight}{HTML}{D4F1F4}

\newcounter{drawex}
\newtcolorbox{drawex}[1]{
  enhanced,breakable,
  colback=drawteallight,colframe=drawteal!70,
  title={\bfseries Draw on Paper \refstepcounter{drawex}\thedrawex: #1},
  fonttitle=\color{white},
  attach boxed title to top left={xshift=6pt,yshift=-3pt},
  boxed title style={colback=drawteal,boxrule=0pt,arc=2pt},
  arc=2pt,boxrule=0.4pt,
  top=12pt,
}


\newcommand{\takehomebanner}{%
  \vspace{6pt}
  \begin{tcolorbox}[
    colback=warn!10,colframe=warn,boxrule=0.5pt,arc=2pt,
    left=10pt,right=10pt,top=6pt,bottom=6pt
  ]
    \textbf{Take-Home (Paper Work).}\;
    Solve the following on paper without using a computer or compiler.
    Hand-write your answers; submit at the start of the next lab.
    Goal: prepare you to dry-code during the written exam.
  \end{tcolorbox}
  \vspace{4pt}
}
 after \documentclass{article}.

\usepackage[a4paper,margin=2.2cm]{geometry}
\usepackage[T1]{fontenc}
\usepackage{lmodern}
\usepackage{microtype}
\usepackage{enumitem}
\usepackage{xcolor}
\usepackage{tcolorbox}
\tcbuselibrary{breakable,skins}
\usepackage{listings}
\usepackage{fancyhdr}
\usepackage{titlesec}
\usepackage{amsmath}
\usepackage{array}
\usepackage{booktabs}
\usepackage{hyperref}

\hypersetup{
  colorlinks=true,
  linkcolor=black,
  urlcolor=blue!60!black,
  pdfborder={0 0 0}
}

% ---------- Course identity (override per file if needed) ----------
\providecommand{\CourseName}{Microprocessors}
\providecommand{\CourseCode}{010153521}
\providecommand{\Semester}{Semester 1/2026}
\providecommand{\LabDuration}{3 hours (in-lab)}

% ---------- Colors ----------
\definecolor{accent}{HTML}{1F4E79}
\definecolor{accentlight}{HTML}{D9E2EC}
\definecolor{warn}{HTML}{B45309}
\definecolor{ok}{HTML}{166534}
\definecolor{codebg}{HTML}{F6F8FA}
\definecolor{codeframe}{HTML}{D0D7DE}
\definecolor{codekw}{HTML}{0550AE}
\definecolor{codestr}{HTML}{0A7B2C}
\definecolor{codecom}{HTML}{6E7781}

% ---------- Code / assembly listings ----------
\lstdefinestyle{c}{
  language=C,
  basicstyle=\ttfamily\small,
  keywordstyle=\color{codekw}\bfseries,
  stringstyle=\color{codestr},
  commentstyle=\color{codecom}\itshape,
  backgroundcolor=\color{codebg},
  frame=single,
  rulecolor=\color{codeframe},
  framerule=0.4pt,
  showstringspaces=false,
  breaklines=true,
  columns=fullflexible,
  keepspaces=true,
  tabsize=4,
  xleftmargin=8pt,
  xrightmargin=4pt,
  aboveskip=6pt,
  belowskip=6pt,
}
\lstdefinestyle{asm}{
  basicstyle=\ttfamily\small,
  commentstyle=\color{codecom}\itshape,
  backgroundcolor=\color{codebg},
  frame=single,
  rulecolor=\color{codeframe},
  framerule=0.4pt,
  showstringspaces=false,
  breaklines=true,
  columns=fullflexible,
  keepspaces=true,
  tabsize=4,
  xleftmargin=8pt,
  xrightmargin=4pt,
  aboveskip=6pt,
  belowskip=6pt,
}
\lstset{style=asm}

% ---------- Section styling ----------
\titleformat{\section}{\Large\bfseries\color{accent}}{\thesection.}{0.5em}{}
\titleformat{\subsection}{\large\bfseries\color{accent!85!black}}{\thesubsection}{0.5em}{}
\titlespacing*{\section}{0pt}{12pt}{6pt}
\titlespacing*{\subsection}{0pt}{8pt}{4pt}

% ---------- Page headers/footers ----------
\pagestyle{fancy}
\fancyhf{}
\fancyhead[L]{\small\CourseName\ifx\CourseCode\empty\else\ (\CourseCode)\fi}
\fancyhead[C]{\small\Semester}
\fancyhead[R]{\small\textbf{Lab \LabNumber:} \LabTitle}
\fancyfoot[C]{\small\thepage}
\renewcommand{\headrulewidth}{0.4pt}

% ---------- Title block ----------
\newcommand{\labheader}{%
  \begin{tcolorbox}[
    enhanced,
    colback=accentlight,
    colframe=accent,
    boxrule=0.5pt,
    arc=2pt,
    left=10pt, right=10pt, top=8pt, bottom=8pt
  ]
    {\Large\bfseries\color{accent} Lab \LabNumber: \LabTitle}\\[2pt]
    {\small \CourseName\ifx\CourseCode\empty\else\ (\CourseCode)\fi\quad\textbar\quad \Semester}\\[1pt]
    {\small \textbf{Duration:} \LabDuration\quad\textbar\quad \textbf{Topic:} \LabTopic}
  \end{tcolorbox}
  \vspace{2pt}
}

% ---------- Custom environments ----------
\newtcolorbox{summary}[1][Quick Reference]{
  enhanced,breakable,
  colback=accentlight!40,colframe=accent!70,
  title={\bfseries #1},
  fonttitle=\color{white},
  attach boxed title to top left={xshift=6pt,yshift=-3pt},
  boxed title style={colback=accent,boxrule=0pt,arc=2pt},
  arc=2pt,boxrule=0.4pt,
  top=12pt,
}

\newcounter{labex}[section]
\newtcolorbox{labexercise}[1]{
  enhanced,breakable,
  colback=white,colframe=accent,
  title={\bfseries In-Lab Exercise \refstepcounter{labex}\thelabex: #1},
  fonttitle=\color{white},
  attach boxed title to top left={xshift=6pt,yshift=-3pt},
  boxed title style={colback=accent,boxrule=0pt,arc=2pt},
  arc=2pt,boxrule=0.4pt,
  top=12pt,
}

\newcounter{traceex}
\newtcolorbox{trace}[1][]{
  enhanced,breakable,
  colback=white,colframe=warn,
  title={\bfseries Trace \& Predict \refstepcounter{traceex}\thetraceex\ifx\\#1\\\else: #1\fi},
  fonttitle=\color{white},
  attach boxed title to top left={xshift=6pt,yshift=-3pt},
  boxed title style={colback=warn,boxrule=0pt,arc=2pt},
  arc=2pt,boxrule=0.4pt,
  top=12pt,
}

\newcounter{bugex}
\newtcolorbox{bug}[1][]{
  enhanced,breakable,
  colback=white,colframe=warn,
  title={\bfseries Find the Bug \refstepcounter{bugex}\thebugex\ifx\\#1\\\else: #1\fi},
  fonttitle=\color{white},
  attach boxed title to top left={xshift=6pt,yshift=-3pt},
  boxed title style={colback=warn,boxrule=0pt,arc=2pt},
  arc=2pt,boxrule=0.4pt,
  top=12pt,
}

\newcounter{writeex}
\newtcolorbox{writecode}[1][]{
  enhanced,breakable,
  colback=white,colframe=warn,
  title={\bfseries Write Code on Paper \refstepcounter{writeex}\thewriteex\ifx\\#1\\\else: #1\fi},
  fonttitle=\color{white},
  attach boxed title to top left={xshift=6pt,yshift=-3pt},
  boxed title style={colback=warn,boxrule=0pt,arc=2pt},
  arc=2pt,boxrule=0.4pt,
  top=12pt,
}

% ---------- AI Collaboration box ----------
\definecolor{aipurple}{HTML}{4C1D95}
\definecolor{aipurplelight}{HTML}{EDE9FE}

\newtcolorbox{aicollab}[1][AI Collaboration]{
  enhanced,breakable,
  colback=aipurplelight,colframe=aipurple!70,
  title={\bfseries #1},
  fonttitle=\color{white},
  attach boxed title to top left={xshift=6pt,yshift=-3pt},
  boxed title style={colback=aipurple,boxrule=0pt,arc=2pt},
  arc=2pt,boxrule=0.4pt,
  top=12pt,
}

\newcounter{tryex}
\newtcolorbox{tryit}[1]{
  enhanced,breakable,
  colback=ok!6,colframe=ok!70!black,
  title={\bfseries Try It \refstepcounter{tryex}\thetryex: #1},
  fonttitle=\color{white},
  attach boxed title to top left={xshift=6pt,yshift=-3pt},
  boxed title style={colback=ok!70!black,boxrule=0pt,arc=2pt},
  arc=2pt,boxrule=0.4pt,
  top=12pt,
}

% ---------- Drawing exercise ----------
\definecolor{drawteal}{HTML}{0D7377}
\definecolor{drawteallight}{HTML}{D4F1F4}

\newcounter{drawex}
\newtcolorbox{drawex}[1]{
  enhanced,breakable,
  colback=drawteallight,colframe=drawteal!70,
  title={\bfseries Draw on Paper \refstepcounter{drawex}\thedrawex: #1},
  fonttitle=\color{white},
  attach boxed title to top left={xshift=6pt,yshift=-3pt},
  boxed title style={colback=drawteal,boxrule=0pt,arc=2pt},
  arc=2pt,boxrule=0.4pt,
  top=12pt,
}


\newcommand{\takehomebanner}{%
  \vspace{6pt}
  \begin{tcolorbox}[
    colback=warn!10,colframe=warn,boxrule=0.5pt,arc=2pt,
    left=10pt,right=10pt,top=6pt,bottom=6pt
  ]
    \textbf{Take-Home (Paper Work).}\;
    Solve the following on paper without using a computer or compiler.
    Hand-write your answers; submit at the start of the next lab.
    Goal: prepare you to dry-code during the written exam.
  \end{tcolorbox}
  \vspace{4pt}
}
 after \documentclass{article}.

\usepackage[a4paper,margin=2.2cm]{geometry}
\usepackage[T1]{fontenc}
\usepackage{lmodern}
\usepackage{microtype}
\usepackage{enumitem}
\usepackage{xcolor}
\usepackage{tcolorbox}
\tcbuselibrary{breakable,skins}
\usepackage{listings}
\usepackage{fancyhdr}
\usepackage{titlesec}
\usepackage{amsmath}
\usepackage{array}
\usepackage{booktabs}
\usepackage{hyperref}

\hypersetup{
  colorlinks=true,
  linkcolor=black,
  urlcolor=blue!60!black,
  pdfborder={0 0 0}
}

% ---------- Course identity (override per file if needed) ----------
\providecommand{\CourseName}{Microprocessors}
\providecommand{\CourseCode}{010153521}
\providecommand{\Semester}{Semester 1/2026}
\providecommand{\LabDuration}{3 hours (in-lab)}

% ---------- Colors ----------
\definecolor{accent}{HTML}{1F4E79}
\definecolor{accentlight}{HTML}{D9E2EC}
\definecolor{warn}{HTML}{B45309}
\definecolor{ok}{HTML}{166534}
\definecolor{codebg}{HTML}{F6F8FA}
\definecolor{codeframe}{HTML}{D0D7DE}
\definecolor{codekw}{HTML}{0550AE}
\definecolor{codestr}{HTML}{0A7B2C}
\definecolor{codecom}{HTML}{6E7781}

% ---------- Code / assembly listings ----------
\lstdefinestyle{c}{
  language=C,
  basicstyle=\ttfamily\small,
  keywordstyle=\color{codekw}\bfseries,
  stringstyle=\color{codestr},
  commentstyle=\color{codecom}\itshape,
  backgroundcolor=\color{codebg},
  frame=single,
  rulecolor=\color{codeframe},
  framerule=0.4pt,
  showstringspaces=false,
  breaklines=true,
  columns=fullflexible,
  keepspaces=true,
  tabsize=4,
  xleftmargin=8pt,
  xrightmargin=4pt,
  aboveskip=6pt,
  belowskip=6pt,
}
\lstdefinestyle{asm}{
  basicstyle=\ttfamily\small,
  commentstyle=\color{codecom}\itshape,
  backgroundcolor=\color{codebg},
  frame=single,
  rulecolor=\color{codeframe},
  framerule=0.4pt,
  showstringspaces=false,
  breaklines=true,
  columns=fullflexible,
  keepspaces=true,
  tabsize=4,
  xleftmargin=8pt,
  xrightmargin=4pt,
  aboveskip=6pt,
  belowskip=6pt,
}
\lstset{style=asm}

% ---------- Section styling ----------
\titleformat{\section}{\Large\bfseries\color{accent}}{\thesection.}{0.5em}{}
\titleformat{\subsection}{\large\bfseries\color{accent!85!black}}{\thesubsection}{0.5em}{}
\titlespacing*{\section}{0pt}{12pt}{6pt}
\titlespacing*{\subsection}{0pt}{8pt}{4pt}

% ---------- Page headers/footers ----------
\pagestyle{fancy}
\fancyhf{}
\fancyhead[L]{\small\CourseName\ifx\CourseCode\empty\else\ (\CourseCode)\fi}
\fancyhead[C]{\small\Semester}
\fancyhead[R]{\small\textbf{Lab \LabNumber:} \LabTitle}
\fancyfoot[C]{\small\thepage}
\renewcommand{\headrulewidth}{0.4pt}

% ---------- Title block ----------
\newcommand{\labheader}{%
  \begin{tcolorbox}[
    enhanced,
    colback=accentlight,
    colframe=accent,
    boxrule=0.5pt,
    arc=2pt,
    left=10pt, right=10pt, top=8pt, bottom=8pt
  ]
    {\Large\bfseries\color{accent} Lab \LabNumber: \LabTitle}\\[2pt]
    {\small \CourseName\ifx\CourseCode\empty\else\ (\CourseCode)\fi\quad\textbar\quad \Semester}\\[1pt]
    {\small \textbf{Duration:} \LabDuration\quad\textbar\quad \textbf{Topic:} \LabTopic}
  \end{tcolorbox}
  \vspace{2pt}
}

% ---------- Custom environments ----------
\newtcolorbox{summary}[1][Quick Reference]{
  enhanced,breakable,
  colback=accentlight!40,colframe=accent!70,
  title={\bfseries #1},
  fonttitle=\color{white},
  attach boxed title to top left={xshift=6pt,yshift=-3pt},
  boxed title style={colback=accent,boxrule=0pt,arc=2pt},
  arc=2pt,boxrule=0.4pt,
  top=12pt,
}

\newcounter{labex}[section]
\newtcolorbox{labexercise}[1]{
  enhanced,breakable,
  colback=white,colframe=accent,
  title={\bfseries In-Lab Exercise \refstepcounter{labex}\thelabex: #1},
  fonttitle=\color{white},
  attach boxed title to top left={xshift=6pt,yshift=-3pt},
  boxed title style={colback=accent,boxrule=0pt,arc=2pt},
  arc=2pt,boxrule=0.4pt,
  top=12pt,
}

\newcounter{traceex}
\newtcolorbox{trace}[1][]{
  enhanced,breakable,
  colback=white,colframe=warn,
  title={\bfseries Trace \& Predict \refstepcounter{traceex}\thetraceex\ifx\\#1\\\else: #1\fi},
  fonttitle=\color{white},
  attach boxed title to top left={xshift=6pt,yshift=-3pt},
  boxed title style={colback=warn,boxrule=0pt,arc=2pt},
  arc=2pt,boxrule=0.4pt,
  top=12pt,
}

\newcounter{bugex}
\newtcolorbox{bug}[1][]{
  enhanced,breakable,
  colback=white,colframe=warn,
  title={\bfseries Find the Bug \refstepcounter{bugex}\thebugex\ifx\\#1\\\else: #1\fi},
  fonttitle=\color{white},
  attach boxed title to top left={xshift=6pt,yshift=-3pt},
  boxed title style={colback=warn,boxrule=0pt,arc=2pt},
  arc=2pt,boxrule=0.4pt,
  top=12pt,
}

\newcounter{writeex}
\newtcolorbox{writecode}[1][]{
  enhanced,breakable,
  colback=white,colframe=warn,
  title={\bfseries Write Code on Paper \refstepcounter{writeex}\thewriteex\ifx\\#1\\\else: #1\fi},
  fonttitle=\color{white},
  attach boxed title to top left={xshift=6pt,yshift=-3pt},
  boxed title style={colback=warn,boxrule=0pt,arc=2pt},
  arc=2pt,boxrule=0.4pt,
  top=12pt,
}

% ---------- AI Collaboration box ----------
\definecolor{aipurple}{HTML}{4C1D95}
\definecolor{aipurplelight}{HTML}{EDE9FE}

\newtcolorbox{aicollab}[1][AI Collaboration]{
  enhanced,breakable,
  colback=aipurplelight,colframe=aipurple!70,
  title={\bfseries #1},
  fonttitle=\color{white},
  attach boxed title to top left={xshift=6pt,yshift=-3pt},
  boxed title style={colback=aipurple,boxrule=0pt,arc=2pt},
  arc=2pt,boxrule=0.4pt,
  top=12pt,
}

\newcounter{tryex}
\newtcolorbox{tryit}[1]{
  enhanced,breakable,
  colback=ok!6,colframe=ok!70!black,
  title={\bfseries Try It \refstepcounter{tryex}\thetryex: #1},
  fonttitle=\color{white},
  attach boxed title to top left={xshift=6pt,yshift=-3pt},
  boxed title style={colback=ok!70!black,boxrule=0pt,arc=2pt},
  arc=2pt,boxrule=0.4pt,
  top=12pt,
}

% ---------- Drawing exercise ----------
\definecolor{drawteal}{HTML}{0D7377}
\definecolor{drawteallight}{HTML}{D4F1F4}

\newcounter{drawex}
\newtcolorbox{drawex}[1]{
  enhanced,breakable,
  colback=drawteallight,colframe=drawteal!70,
  title={\bfseries Draw on Paper \refstepcounter{drawex}\thedrawex: #1},
  fonttitle=\color{white},
  attach boxed title to top left={xshift=6pt,yshift=-3pt},
  boxed title style={colback=drawteal,boxrule=0pt,arc=2pt},
  arc=2pt,boxrule=0.4pt,
  top=12pt,
}


\newcommand{\takehomebanner}{%
  \vspace{6pt}
  \begin{tcolorbox}[
    colback=warn!10,colframe=warn,boxrule=0.5pt,arc=2pt,
    left=10pt,right=10pt,top=6pt,bottom=6pt
  ]
    \textbf{Take-Home (Paper Work).}\;
    Solve the following on paper without using a computer or compiler.
    Hand-write your answers; submit at the start of the next lab.
    Goal: prepare you to dry-code during the written exam.
  \end{tcolorbox}
  \vspace{4pt}
}

\usepackage{tikz}
\usetikzlibrary{positioning,shapes.geometric,arrows.meta}
\usepackage{circuitikz}

\begin{document}
\labheader

\section*{Objectives}
\begin{itemize}[leftmargin=*,nosep]
  \item Design a general I/O circuit that fits inside the Arduino Uno's pin budget by
        \textbf{time-division multiplexing} a 2-digit 7-segment display, instead of
        wiring one display driver pin per segment per digit.
  \item Wire push-button inputs with the ATmega328P's \textbf{internal pull-up}
        resistors (Lecture 4, \S4.2-4.3) and individually addressable LED outputs
        with \texttt{SBI}/\texttt{CBI} (Lecture 4, \S4.4-4.5) on a breadboard, driven
        by the same Arduino Uno R3 toolchain as Lab 3.
  \item Write and test AVR assembly that scans buttons with \texttt{SBIC}/\texttt{SBIS},
        converts a digit to a 7-segment pattern, and refreshes a multiplexed display
        fast enough to avoid visible flicker -- applying Lecture 3's cycle-counting
        delay methods to a real refresh-rate budget.
  \item Produce a verified, working breadboard circuit and firmware -- this is the
        exact circuit you will capture into a KiCad schematic and PCB in
        \textbf{Lab 4-2}.
\end{itemize}

\begin{summary}[Quick Reference: Pin Assignment]
Every signal fits inside the 18 pins Lab 3 already established as usable on the
Arduino Uno R3 (ATmega328P) -- \texttt{PB0-PB5}, \texttt{PC0-PC5}, and
\texttt{PD2-PD7} (\texttt{PD0}/\texttt{PD1} stay off-limits, reserved for USB-serial).
\begin{center}
\begin{tabular}{@{}p{2.0cm}p{1.8cm}p{2.2cm}p{7.9cm}@{}}
\toprule
\textbf{Port bits} & \textbf{Uno pins} & \textbf{Direction} & \textbf{Function} \\
\midrule
\texttt{PB0-PB3} & D8-D11 & Input, pull-up & \texttt{SW1-SW4} push buttons (active
\textbf{low}: pressed = 0) \\
\texttt{PB4}, \texttt{PB5} & D12, D13 & Output & \texttt{DIG1}, \texttt{DIG2} digit-select
transistor bases (D13 doubles as the Uno's on-board LED -- see Try It 1) \\
\texttt{PC0-PC5} & A0-A5 & Output & Segments \texttt{a,b,c,d,e,f} (shared by both
digits) \\
\texttt{PD2} & D2 & Output & Segment \texttt{g} (shared by both digits) \\
\texttt{PD3-PD7} & D3-D7 & Output & \texttt{LED1-LED5}, individually addressable \\
\bottomrule
\end{tabular}
\end{center}
Segments \texttt{a-f} land on \texttt{PORTC} bits \texttt{0-5} in the same bit order
used by the standard segment code table below, so one \texttt{OUT PORTC,Rd} sets all
six at once -- only segment \texttt{g} (on \texttt{PORTD}) needs a separate
\texttt{SBI}/\texttt{CBI}.
\end{summary}

\begin{summary}[Quick Reference: 7-Segment Codes (Common Cathode)]
Both digits are common-\textbf{cathode}: a segment lights when its drive pin is
\textbf{HIGH} (current flows pin $\to$ resistor $\to$ segment LED $\to$ common
cathode). Bit \texttt{0} = segment \texttt{a} \ldots bit \texttt{5} = segment
\texttt{f}, bit \texttt{6} = segment \texttt{g} (bit 7 unused) -- exactly the
classic \texttt{gfedcba} encoding.
\begin{center}
\begin{tabular}{@{}cccccccccc@{}}
\toprule
\textbf{Digit} & 0 & 1 & 2 & 3 & 4 & 5 & 6 & 7 & 8 \\
\midrule
\textbf{Hex} & \texttt{0x3F} & \texttt{0x06} & \texttt{0x5B} & \texttt{0x4F} &
\texttt{0x66} & \texttt{0x6D} & \texttt{0x7D} & \texttt{0x07} & \texttt{0x7F} \\
\bottomrule
\end{tabular}
\end{center}
\textbf{Digit 9} = \texttt{0x6F}. Bits \texttt{0-5} of each byte go straight to
\texttt{PORTC}; bit \texttt{6} (segment \texttt{g}) must be tested separately (e.g.\
\texttt{SBRC}/\texttt{SBRS}) and mirrored onto \texttt{PD2} with \texttt{SBI}/\texttt{CBI}.
\end{summary}

\begin{summary}[Quick Reference: Multiplexing and Refresh Rate]
Only \emph{one} digit's common cathode is ever grounded at a time -- the other
digit's segments have no return path, so only the grounded digit actually lights,
even though both digits' segment pins are wired in parallel. Show each digit for a
short \textbf{dwell time} $t_{\text{dwell}}$, then switch to the other -- fast
enough that the human eye's flicker-fusion threshold ($\approx$50-60~Hz) blends the
two into one steady image:
\[
f_{\text{refresh}} = \frac{1}{N_{\text{digits}} \times t_{\text{dwell}}}, \qquad
N_{\text{digits}} = 2 \ \Rightarrow\ t_{\text{dwell}} \le \frac{1}{2 \times 60\text{Hz}}
\approx 8.3\text{ ms}
\]
Stay comfortably under this bound -- a $t_{\text{dwell}}$ near the limit will show
visible flicker on real hardware even though it looks fine calculated on paper.
\end{summary}

\begin{summary}[Quick Reference: Your Individual Target]
As in Lab 3, derive your personal targets from your student ID before touching the
keyboard:
\[
D = \text{sum of every digit in your student ID}
\]
\begin{center}
\begin{tabular}{@{}p{5.7cm}p{8.9cm}@{}}
\toprule
\textbf{Exercise} & \textbf{Your target} \\
\midrule
3 (multiplexed display) & $t_{\text{dwell}} = 2\text{ ms} + (D \bmod 8) \times 0.5\text{ ms}$ \\
4 (button counter) & $N_0 = D \bmod 100$ (power-on displayed value) \\
\bottomrule
\end{tabular}
\end{center}
\textbf{Worked example} (a made-up ID, do not reuse): ID \texttt{1029384756123}
$\to D = 51$. Then $t_{\text{dwell}} = 2 + (51 \bmod 8) \times 0.5 = 2 + 3\times0.5 =
3.5$~ms (refresh rate $\approx$143 Hz, well clear of the 60 Hz bound), and
$N_0 = 51 \bmod 100 = 51$.
\end{summary}

\section*{Bill of Materials}
\begin{center}
\begin{tabular}{@{}p{1.4cm}p{6.3cm}p{7.5cm}@{}}
\toprule
\textbf{Qty} & \textbf{Part} & \textbf{Notes} \\
\midrule
1 & Arduino Uno R3 (from Lab 3) & Supplies power, ground, and all 18 signal pins \\
4 & Tactile push button (6~mm THT) & No external resistor -- internal pull-ups only \\
5 & LED, 3~mm or 5~mm & Any color; one per \texttt{LED1-LED5} \\
5 & Resistor, 330~$\Omega$ & Current limiting for \texttt{LED1-LED5} \\
2 & 7-segment display, \textbf{common cathode} & \texttt{DIGIT1} (tens), \texttt{DIGIT2} (units) \\
7 & Resistor, 330~$\Omega$ & Shared segment resistors \texttt{a-g} (one per segment, not per digit) \\
2 & NPN transistor (e.g.\ 2N3904) & Digit-select low-side switches \\
2 & Resistor, 1~k$\Omega$ & Transistor base resistors \\
1 & Breadboard + jumper wires & \\
\bottomrule
\end{tabular}
\end{center}

\section*{Building the Circuit}
\begin{enumerate}[leftmargin=*,nosep]
  \item \textbf{Buttons.} Wire one leg of each \texttt{SW1-SW4} to \texttt{PB0-PB3}
        (Uno D8-D11), the other leg to \textbf{GND}. Do \emph{not} add an external
        pull-up or pull-down -- firmware will enable the AVR's internal pull-up, so a
        pressed button reads \texttt{0} and a released one reads \texttt{1}.
  \item \textbf{LEDs.} Wire each of \texttt{LED1-LED5} anode $\to$ 330~$\Omega$
        resistor $\to$ \texttt{PD3-PD7} (Uno D3-D7); cathode $\to$ \textbf{GND}.
  \item \textbf{Segment bus.} For each of the 7 segments (\texttt{a-g}), tie
        \emph{both} digits' same-lettered segment pin together, then run that shared
        node through one 330~$\Omega$ resistor to the matching MCU pin
        (\texttt{a-f} $\to$ \texttt{PC0-PC5}; \texttt{g} $\to$ \texttt{PD2}). This is
        the heart of the multiplexing trick: 7 resistors and 7 MCU pins drive
        \emph{both} digits, not 14.
  \item \textbf{Digit-select transistors.} \texttt{DIGIT1}'s common-cathode pin goes
        to Q1's \textbf{collector}; Q1's \textbf{emitter} goes to GND; Q1's
        \textbf{base} goes through a 1~k$\Omega$ resistor to \texttt{PB4} (D12).
        Wire \texttt{DIGIT2} to Q2 the same way, based from \texttt{PB5} (D13).
        Driving \texttt{PB4}/\texttt{PB5} HIGH saturates the transistor, sinking that
        digit's entire segment current to ground -- the AVR pin only supplies a few
        mA of \emph{base} current, not the segments' combined current.
\end{enumerate}

\begin{figure}[h!]
\centering
\begin{circuitikz}[
  american voltages,
  scale=0.78, transform shape,
  font=\scriptsize
]

% ---------------- Buttons: SW1-SW4 ----------------
\draw (0,10)   node[above]{\texttt{PB0} (D8)}  to[push button, l=SW1] (0,8.5)   node[ground]{};
\draw (1.8,10) node[above]{\texttt{PB1} (D9)}  to[push button, l=SW2] (1.8,8.5) node[ground]{};
\draw (3.6,10) node[above]{\texttt{PB2} (D10)} to[push button, l=SW3] (3.6,8.5) node[ground]{};
\draw (5.4,10) node[above]{\texttt{PB3} (D11)} to[push button, l=SW4] (5.4,8.5) node[ground]{};

% ---------------- LEDs: LED1-LED5 ----------------
\draw (8,10)    node[above]{\texttt{PD3} (D3)} to[R, l=$330\Omega$] (8,8.8)    to[led, l=LED1] (8,7.3)    node[ground]{};
\draw (9.8,10)  node[above]{\texttt{PD4} (D4)} to[R, l=$330\Omega$] (9.8,8.8)  to[led, l=LED2] (9.8,7.3)  node[ground]{};
\draw (11.6,10) node[above]{\texttt{PD5} (D5)} to[R, l=$330\Omega$] (11.6,8.8) to[led, l=LED3] (11.6,7.3) node[ground]{};
\draw (13.4,10) node[above]{\texttt{PD6} (D6)} to[R, l=$330\Omega$] (13.4,8.8) to[led, l=LED4] (13.4,7.3) node[ground]{};
\draw (15.2,10) node[above]{\texttt{PD7} (D7)} to[R, l=$330\Omega$] (15.2,8.8) to[led, l=LED5] (15.2,7.3) node[ground]{};

% ---------------- Representative segment channel (segment a) ----------------
\draw (0,6.5) node[above, align=center]{\texttt{PC0} (A0)\\segment \texttt{a}}
      to[R, l=$330\Omega$] (0,5) coordinate(nodeA);
\node[circ] at (nodeA) {};

\node[draw, fill=accentlight, minimum width=3.2cm, minimum height=1.6cm, align=center]
     at (3,5)   (dig1) {\textbf{DIGIT1}\\segment \texttt{a} (shown)\\\texttt{b--g}: identical};
\node[draw, fill=accentlight, minimum width=3.2cm, minimum height=1.6cm, align=center]
     at (3,1.0) (dig2) {\textbf{DIGIT2}\\segment \texttt{a} (shown)\\\texttt{b--g}: identical};

\draw (nodeA) -| (dig1.west);
\draw (nodeA) |- (dig2.west);

% ---------------- Digit-select transistors ----------------
\node[npn] at (6.5,3) (Q1) {};
\node[npn] at (6.5,-1.4) (Q2) {};

\draw (dig1.south) |- (Q1.C) node[pos=0.9, above right]{CC1};
\draw (dig2.south) |- (Q2.C) node[pos=0.9, above right]{CC2};

\draw (Q1.E) -- ++(0,-0.8) node[ground]{};
\draw (Q2.E) -- ++(0,-0.8) node[ground]{};

\draw (Q1.B) -- ++(-1.2,0) to[R, l=1k$\Omega$] ++(-1.6,0) node[left]{\texttt{PB4} (D12)};
\draw (Q2.B) -- ++(-1.2,0) to[R, l=1k$\Omega$] ++(-1.6,0) node[left]{\texttt{PB5} (D13)};

\end{circuitikz}
\caption{Full schematic for buttons (\texttt{SW1-SW4}) and LEDs (\texttt{LED1-LED5}),
plus the multiplexed-display driver with segment \texttt{a} drawn in full -- segments
\texttt{b-g} repeat the identical pattern (own 330~$\Omega$ resistor from
\texttt{PC1-PC5}/\texttt{PD2}, tied to both digits' matching pin) and are omitted here
only to keep the figure readable. \texttt{CC1}/\texttt{CC2} are each digit's common
cathode; only one of Q1/Q2 conducts at a time, so only one digit's segments have a
path to ground at any instant.}
\end{figure}

\section*{Quick Experiments (Try It)}

\begin{tryit}{The on-board LED is watching PB5}
Before wiring anything else, load Lab 3's blink skeleton (or any program that
toggles \texttt{PB5}) and confirm the Uno's on-board D13 LED blinks. Now wire
\texttt{DIG2} (Exercise 3 below) to \texttt{PB5} as instructed above and re-run your
multiplexed display firmware.
\begin{enumerate}[nosep,leftmargin=*]
  \item Does the on-board LED flicker at the same rate as \texttt{DIGIT2}'s refresh?
  \item Explain why sharing a pin between the on-board LED and \texttt{DIG2}'s
        transistor base does \emph{not} stop the display from working -- what does
        the on-board LED's own series resistor mean for how much current it steals
        from the transistor's base?
\end{enumerate}
\end{tryit}

\begin{tryit}{One digit, no multiplexing}
Wire \emph{only} \texttt{DIGIT1} (leave \texttt{DIGIT2} unpowered, \texttt{PB5} low).
Write a short loop that walks digit values 0 through 9, holding each for about half a
second, using the segment code table above and a single \texttt{OUT PORTC} +
\texttt{SBI}/\texttt{CBI PORTD,2} per digit. Confirm every digit 0-9 renders exactly
as expected before moving on to multiplexing two digits at once -- debugging one
static digit is far easier than debugging two flickering ones.
\end{tryit}

\section*{In-Lab Exercises (target: 3 hours)}

\begin{labexercise}{Buttons and LEDs \hfill {\small \itshape (35 min)}}
Wire \texttt{SW1-SW4} and \texttt{LED1-LED5} per the instructions above. Write and
run a program that mirrors each button directly to one LED (e.g.\ \texttt{SW1}
$\to$ \texttt{LED1}), using the \texttt{SBIC}/\texttt{SBIS} polling pattern from
Lecture 4's ``Reading a Single Bit: Switch to LED'' example.
\begin{itemize}[nosep,leftmargin=*]
  \item Confirm \texttt{DDRB} bits 0-3 are \textbf{0} (input) and \texttt{PORTB}
        bits 0-3 are \textbf{1} (pull-up enabled) -- verify with \texttt{IN}/inspect
        in the simulator or by reasoning through your \texttt{.asm}, not just by
        observing the LED behave correctly (a swapped polarity can look right by
        accident on some button wiring).
  \item Record, in your notebook, what \texttt{LED1} does the instant power is
        applied, \emph{before} you press anything -- and explain why, in terms of
        the pull-up's default state.
\end{itemize}
\end{labexercise}

\begin{labexercise}{Static digit and the segment table \hfill {\small \itshape (30 min)}}
Repeat Try It 2 as a graded checkpoint: \texttt{DIGIT1} alone, cycling 0-9 with a
visible hold time per digit. Write a \texttt{DIGIT\_TO\_SEG} subroutine (input: a
digit 0-9 in a register; output: the segment byte in another register) using a
\texttt{CPI}/\texttt{BRNE} comparison chain -- \emph{not} a lookup table -- so each
comparison is explicit and traceable in your \texttt{.lst} file.
\end{labexercise}

\begin{labexercise}{Multiplexed 2-digit counter \hfill {\small \itshape (55 min)}}
Using your $t_{\text{dwell}}$ from the ``Your Individual Target'' box, build the
full multiplexed \texttt{REFRESH} loop: show the tens digit for $t_{\text{dwell}}$
with \texttt{DIG1} selected, then the units digit for $t_{\text{dwell}}$ with
\texttt{DIG2} selected, repeating forever. Drive a free-running two-digit count
(00-99, wrapping) that increments roughly once per second, independent of the
display refresh loop.
\begin{enumerate}[nosep,leftmargin=*]
  \item Show your cycle accounting for $t_{\text{dwell}}$ (the delay subroutine
        tuned to your target, same method as Lab 3) and your calculated
        $f_{\text{refresh}}$.
  \item Build it. Does the display look like one steady image, or can you see it
        flicker or ``strobe'' when you wave a hand or pencil quickly in front of it?
        If you see flicker, is your measured behavior consistent with your
        calculated $f_{\text{refresh}}$ against the 60~Hz bound?
  \item \textbf{Trace/predict first, then verify:} before running it, predict what
        the display shows if you accidentally swap the \texttt{DIG1}/\texttt{DIG2}
        select lines but leave the tens/units segment data unswapped. Then swap them
        on purpose and confirm.
\end{enumerate}
\end{labexercise}

\begin{labexercise}{Button-controlled counter \hfill {\small \itshape (40 min)}}
Using your $N_0$ from the ``Your Individual Target'' box, initialize the displayed
value to $N_0$ at power-on. Wire \texttt{SW1} to increment and \texttt{SW2} to
decrement the displayed 2-digit value (wrapping 99$\to$00 and 00$\to$99), each on a
single button press -- not continuously while held.
\begin{itemize}[nosep,leftmargin=*]
  \item A mechanical button ``bounces'' -- its contact opens and closes several
        times over a few milliseconds before settling. Add a simple software debounce
        (a short delay after detecting a press, then re-checking the pin is still
        low) and describe in your notebook what you observed \emph{without} it
        (e.g.\ the count jumping by more than 1 per press).
  \item Confirm the power-on value matches your $N_0$ exactly before testing the
        buttons.
\end{itemize}
\end{labexercise}

\begin{aicollab}
\textbf{Try these prompts with an AI tool:}
\begin{itemize}[leftmargin=*,nosep]
  \item ``I'm multiplexing a 2-digit common-cathode 7-segment display from an AVR.
        What dwell time per digit keeps me safely under the flicker-fusion
        threshold, and how does that change with more digits?''
  \item ``How do I use \texttt{SBIC}/\texttt{SBIS} to poll a push button wired with
        an internal pull-up, and what's a simple software debounce technique in AVR
        assembly?''
  \item ``What's the difference between common-cathode and common-anode 7-segment
        segment codes, and how do I convert one table into the other?''
\end{itemize}

\medskip\textbf{Critical check -- verify, don't just accept:}
If an AI tool gives you a segment code table, check every byte against the Quick
Reference table above bit by bit -- segment-code tables online are inconsistent
about bit order (\texttt{gfedcba} vs.\ \texttt{dp\,g\,f\,e\,d\,c\,b\,a}) and about
common-cathode vs.\ common-anode polarity, and a wrong assumption here will show a
plausible-looking but wrong digit on real hardware.

\medskip\textbf{Know the limit:}
AI tools are prone to computing your $t_{\text{dwell}}$ delay-loop cycle count
without accounting for the \texttt{REFRESH} loop's own overhead -- the
\texttt{DIGIT\_TO\_SEG} comparison chain and the \texttt{SBI}/\texttt{CBI} for
segment \texttt{g} all eat into your dwell-time budget, just like \texttt{RET}
overhead did in Lab 3. A plausible-looking calculation can still miss the real
$f_{\text{refresh}}$ by enough to cause visible flicker -- verify on the real
hardware, not just on paper.
\end{aicollab}

\takehomebanner

\begin{trace}[Refresh rate at a different dwell time]
A different student's design uses a \emph{4-digit} multiplexed display (not this
lab's 2-digit design) with $t_{\text{dwell}} = 3$~ms per digit. Calculate
$f_{\text{refresh}}$ and state whether it is above or below the 60~Hz flicker-fusion
bound used in this lab -- would you expect visible flicker?
\end{trace}

\begin{bug}[Segment bus tied to the wrong port]
A student wired segments \texttt{a-f} to \texttt{PC0-PC5} correctly, but wired
segment \texttt{g} to \texttt{PD3} instead of \texttt{PD2} (an easy off-by-one on a
breadboard) -- and then wrote firmware that (correctly, per the lab's instructions)
sets/clears \texttt{PD2} for segment \texttt{g}. \texttt{DDRD} and \texttt{PORTD}
were otherwise configured exactly as this lab specifies for \texttt{LED1-LED5}.
Predict exactly what appears on the display (both digits, all ten digit values) and
separately what happens to \texttt{LED1}, and explain both in terms of what pin the
firmware is actually toggling versus what pin segment \texttt{g} is actually wired
to.
\end{bug}

\begin{writecode}[Segment code for a common-\emph{anode} display]
On paper, without a computer, simulator, or AI tool, rederive the segment byte for
digits 0, 5, and 8 if \texttt{DIGIT1}/\texttt{DIGIT2} were common-\textbf{anode}
displays instead of common-cathode (segment lights when its drive pin is
\textbf{LOW}, and the ``common'' pin is tied to \textbf{+5V} through the
digit-select transistor instead of to GND). State whether the digit-select
transistor type (NPN vs.\ PNP) would also need to change, and why.
\end{writecode}

\begin{drawex}{Predict, then verify, the multiplexed waveform}
Before building Exercise 3, sketch on paper the voltage you expect at \texttt{PB4}
(\texttt{DIG1} select) and \texttt{PB5} (\texttt{DIG2} select) on the same time axis:
axes labeled, both logic levels marked, and $t_{\text{dwell}}$ and the overall
refresh period $2\times t_{\text{dwell}}$ labeled with your calculated numbers. After
building it, if you have access to an oscilloscope or logic analyzer, capture the
real waveform and annotate your sketch with the measured values next to your
predicted ones.
\end{drawex}

\section*{Reflection}

In one short paragraph, compare this lab's multiplexed display to the naive
alternative: wiring one segment-driver pin per segment per digit (7 segments
$\times$ 2 digits = 14 pins, plus 2 more for the commons -- more pins than the Uno
has spare once the buttons and LEDs are wired). What did multiplexing buy you, and
what did it cost -- in firmware complexity (\texttt{DIGIT\_TO\_SEG}, the
\texttt{REFRESH} loop, the dwell-time budget) or in extra components (two
transistors, seven shared resistors) that a one-pin-per-segment design would not
have needed? This exact circuit becomes a KiCad schematic and PCB in
\textbf{Lab 4-2} -- which of today's breadboard mistakes, if any, do you expect
would be harder to catch once it is a fabricated board instead of a breadboard you
can freely rewire?

\end{document}
